% !PRINTED 32
\weberpage{32}
Our whole problem is accordingly contained in equation (6), which says
nothing other than that the function $f(U,V)$, of degree $4n$, is to
contain $2n-2$ quadratic factors. To satisfy the $(2n-2)$ conditions
that follow from this one has at one's disposal the $2n+2$ coefficients
contained in $U,V$, so that four of these remain undetermined. This was
to be foreseen, since arbitrary homogeneous linear functions of
$\xi,\eta$ may be introduced in place of $\xi,\eta$. These four
supernumerary constants can be used to bring the differential (8) into a
normal form. But since the equations of condition for the coefficients
of $U,V$ are not linear, there are several transformations for a given
degree of transformation.

We shall meet this transformation problem again later from a quite
different side, and shall have to go into it far more deeply. The
details of the algebraic problem will therefore not be entered into more
closely here; instead we wish to illustrate what has been said by two
examples.

\begin{center}
\S\,9. \textbf{\gesperrt{The Transformation of the Second Degree.}}
\end{center}

In order to obtain the elliptic differentials in the normal form we set:
\[
\begin{aligned}
f(x,y) &= xy(x-y)(x-\lambda^2 y),\\
\varphi(\xi,\eta) &= \xi\eta(\xi-\eta)(\xi-\varkappa^2\eta),
\end{aligned}
\]
and let $U, V$ be of the second degree.

Equation (6), \S\,8, now becomes, since $T$ is of the second degree
\weq{1}{UV(U-V)(U-\lambda^2 V) = (a\xi+b\eta)^2(a'\xi+b'\eta)^2\,\xi\eta(\xi-\eta)(\xi-\varkappa^2\eta),}
and two of the four factors of the second degree on the left side of
this equation must be squares. We wish to restrict the great number of
possibilities lying in this still further by assuming that $\eta$ and
$y$ are to vanish simultaneously, so that $V$ is divisible by $\eta$.
This gives (apart from a constant factor) the following three
possibilities for $V$:
\[
1.\ V = \xi\eta, \qquad 2.\ V = \eta(\xi-\eta), \qquad 3.\ V = \eta(\xi-\varkappa^2\eta).
\]
Each of these three cases now again comprises three subcases, since of
the three remaining factors $U$, $U-V$, $U-\lambda^2 V$ any two may be
assumed to be squares. Of these latter
%CONT>>
